% ID: SEL-D0022-OA-E04
% ESTADO: cotejada. Segunda pasada visual completa y borrador compilado revisados.
% Fuente: Selectividad/matematicas_aplicadas/reserva_a_EXAMEN 5 -ANDALUCÍA.pdf
% SHA-256: 668689642ca997cf817c9f9f6608d2835efe898e3162b050d7b22b51f225d363
% Página: 1; opción A; ejercicio original 4.

\begin{enumerate}[label=\alph*)]
\item \textbf{(1.5 puntos)} La talla de los individuos de una población sigue una distribución Normal con desviación típica 8 cm y media desconocida. A partir de una muestra aleatoria se ha obtenido un intervalo de confianza al 95\% para estimar la talla media poblacional, que ha resultado ser $(164.86,171.14)$ en cm.

Calcule la talla media de la muestra y el tamaño muestral mínimo necesario para reducir a la mitad el error máximo de estimación anterior.
\item \textbf{(1 punto)} En un club privado con 243 usuarios se ha seleccionado una muestra para hacer un sondeo, según la actividad realizada y por muestreo aleatorio estratificado. En esa muestra, 5 usuarios practican Yoga, 7 Pilates y 15 Mantenimiento, \textquestiondown{}cuántos usuarios están inscritos en cada actividad en ese club?
\end{enumerate}
